\answer{ex:01:1}{(a)~\nt{GLUT} $8\cdot10^{10}$개는 지구 인구의 $10$배이다. 브로드캐스트 뉴런 $\approx1.9\cdot10^6$개는 빈(Wien) 크기의 도시이다. (b)~말단은 $\approx8\cdot10^{14}$개이고 그중 브로드캐스트 말단이 $\approx3\cdot10^{11}$개, 비율 $\approx4\cdot10^{-4}$로, 이 뉴런들의 비율($2.2\cdot10^{-5}$)보다 $\approx16$배 크다.}
\answer{ex:01:2}{(a)~$\tau_c\ln(c_0/c_*)\approx4.6\ms$. (b)~$T\ge\tau_c\ln101$이면 선로가 비워진다. 초당 $\approx220$번 대신 $\approx22$번까지이다.}
\answer{ex:01:3}{$V(s_k)=\gamma^{K-1-k}r$; $\delta_{\text{lamp}}=\gamma^Kr=re^{-T/\tau_\gamma}\approx0.60\,r$(어떤 $\Delta$에서도), $\tau_\gamma\approx1.95\,\text{s}$.}
\answer{ex:01:4}{$n^*=93$, $47$, $25$, $12$. $\alpha$가 작으면 $n^*\approx5/\alpha$이다.}
\answer{ex:01:5}{\nt{GLUT}: $1\to10\to10^4\to10^7$, 뉴런 $\approx10^7$개, $4$단계, $4\ms$. \nt{DOPA}: $1\to10\to3.2\cdot10^5$, 뉴런 $\approx3\cdot10^5$개, $3$단계로 $30$배 적다. \nt{DOPA} 뉴런 수와 같은 자릿수이다.}
\answer{ex:01:6}{$\lambda=f_EJ_E-f_IJ_I$에서 $J_I=37.5$; 전류의 비 $f_IJ_I/f_EJ_E=0.94$; 억제를 겨우 $6.7\,\%$만 약화해도 사태($\lambda\ge1$)가 일어난다.}
\answer{ex:01:7}{각 상태에서 지연마다 $n\ge57$번 제시해야 한다. 같은 감소는 예컨대 $T$와 함께 커지는 내부 시간 추정의 부정확성으로도 생길 수 있다.}
